\section{Introduction}
% DRAFT (2026-06-17). Prose to own/refine; see docs/problem_statement.md. Keep claims modest.

Power transformers are among the most critical and expensive assets in an electrical
network, and an in-service failure can cause extended outages and high repair costs.
Dissolved Gas Analysis (DGA) is the most widely used technique for the early detection
of incipient faults: thermal and electrical stresses decompose the insulating oil and
paper, releasing characteristic gases---hydrogen, methane, ethane, ethylene, acetylene,
and the carbon oxides---whose relative amounts indicate the type of fault
\cite{ieee_c57104,duval2002review}.

Conventional interpretation methods, namely the IEC~60599 ratios, the Rogers ratios, and
the Duval Triangle, are rule-based. They encode expert knowledge as fixed thresholds,
require an expert to apply and arbitrate, and return a ``not determined'' result for a
non-negligible share of samples \cite{iec60599,duval2002review}. Supervised machine
learning can reach very high accuracy; a recent example fuses the conventional methods
into a single score and feeds an ensemble classifier, reporting $99.17\%$ accuracy
\cite{hybrid_aiscore2026}. Such methods, however, require labelled fault data---and often
the conventional rules themselves as inputs---which utilities rarely have at fleet scale.

This paper studies a label-free alternative on a real operational dataset of 4{,}563 main-tank
samples from 628 transformers (2019--2024), and makes three contributions. We first observe
that representations built on gas concentrations encode overall gassing \emph{magnitude} rather
than fault \emph{type}, and show that the centred log-ratio (CLR) of the combustible-gas
composition separates the two by construction, raising agreement with the Duval fault types
from an adjusted Rand index of $0.16$ to $0.55$ with a simple linear projection; a neural
autoencoder (SD-CAE) and an adversarial disentanglement term add nothing, which we report as
negative ablations. We then show, second, that the
field-event notes routinely used to validate such models are, in this fleet, \emph{confounded by
surveillance bias}: a trivial sample count out-predicts the gas chemistry, and a label-free score
built from each unit's early history loses its apparent forward validity once this count is
controlled---a caution we support with a confound-free,
chemistry-defined evaluation. Third, we build a label-free health/risk screening index that
outranks conventional severity baselines, reported together with this confound floor.

The contributions of this work are:
\begin{itemize}
  \item a label-free \emph{severity-disentangled compositional representation}: the centred
        log-ratio of the combustible-gas composition, on which a linear two-dimensional
        projection recovers Duval fault-type structure (ARI $0.16\rightarrow0.55$); the neural
        (SD-CAE) and adversarial-independence variants are reported as negative ablations;
  \item a finding that, in this fleet, the field-event label used to validate such models is
        confounded by surveillance bias---a trivial sample count out-predicts the gas chemistry
        and the apparent forward validity
        vanishes once it is controlled---together with a confound-free, chemistry-defined
        evaluation target;
  \item a label-free health/risk screening index (severity, acetylene-aware, temporal, and
        anomaly terms) that outranks conventional severity baselines, reported against the
        confound floor.
\end{itemize}

The remainder of the paper is organised as follows. Section~\ref{sec:related} reviews related
work; Section~\ref{sec:data} describes the dataset and preprocessing; Section~\ref{sec:method} the
method; and Sections~\ref{sec:exp}--\ref{sec:results} the experiments and results, before the
conclusion.
